% AC-S47: exact leakage interpolation and the reduced ground-line gate.
% New proof from already audited S30--S33; no new prolate asymptotic input.
\section{Exact leakage interpolation and a smaller ground-line gate}
\label{sec:spectral-prolate-leakage-interpolation}

The pointwise exterior estimate in AC-S31 was used there as an
upper bound for both regularity and size.  Its size bound loses a
factor of the window parameter.  The exact leakage identity supplies
the smaller size norm, and interpolation retains it.  This gives a
stronger estimate for the same original trial, not a replacement
trial or a different cutoff.

Use the notation and Fourier normalization of AC-S30--AC-S33:
\(P_\lambda=1_{[-\lambda,\lambda]}\),
\(\mathcal Ff(v)=\int f(u)e^{-2\pi iuv}\,du\),
\(L=\log\lambda\), and
\(\delta_\lambda=1-\nu_4\).  The mean-zero two-mode raw vector
\(\psi_\lambda\) and its sharply truncated theta image
\(k_\lambda\) are exactly those of (SR.2).  In particular
\[
        \|k_\lambda-K\|_2=O(\lambda^{-1/2}),
        \qquad K\ne0,
        \tag{LI.1}
\]
with the scalar and Fourier transform of \(K\) already fixed in
AC-S04 and AC-S30.  Every function on a window is extended by zero
when a whole-line norm is taken.

\begin{theorem}[AC-S47: leakage-scale Rayleigh bound and reduced selection gate]
\label{AC-S47}
For a fixed even prolate index \(n\), let
\(y_n=\mathcal F\phi_{n,\lambda}\), where
\(\phi_{n,\lambda}\) is unit-norm and time-supported, and put
\[
 w_n(t)=1_{(0,\infty)}(t)\sqrt\lambda\,e^{t/2}
                  y_n(\lambda e^t),
 \qquad \delta_n=1-\nu_n.
 \tag{LI.2}
\]
For every integer \(m\ge0\) and \(0<s<1/4\), unconditionally,
\[
\begin{aligned}
 \|t^m w_n\|_2
       &\le C_{n,m}\sqrt{\delta_n}(1+\log\lambda)^m,\\
 \|t^m w_n\|_{H^s}
       &\le C_{n,m,s}\sqrt{\delta_n}\,
           \lambda^{8s}(1+\log\lambda)^{m(1-4s)}.
\end{aligned}
 \tag{LI.3}
\]
for all sufficiently large \(\lambda\).

Under \(A\), for every \(\varepsilon>0\), the original trial satisfies
\[
 q_L(k_\lambda)\le C_\varepsilon\lambda^\varepsilon
                                  (1-\nu_4),
 \qquad
 a_1(\log\lambda)\le C_\varepsilon
                     \lambda^{9+\varepsilon}e^{-4\pi\lambda^2}.
 \tag{LI.4}
\]
The same form bound holds for the even part of \(k_\lambda\).

Consequently the following strictly weaker replacement of (SE.14)
is sufficient for its full ground-line conclusion: assume \(A\) and
suppose that some \(c>0\) satisfies
\[
 a_2(\log\lambda)\ge
               c\lambda^{10}e^{-4\pi\lambda^2}
              \quad\hbox{for all sufficiently large }\lambda.
 \tag{LI.5}
\]
Then the actual lowest eigenvalue is eventually simple and even.
With its unit eigenvector \(u_\lambda\) phased against the even trial,
\[
 \int_{\mathbb R}e^{b|x|}
       \bigl|\|K\|_2u_\lambda(x)-K(x)\bigr|\,dx
                 \longrightarrow0
                  \qquad(0\le b<1/2).
 \tag{LI.6}
\]
Thus the same scalar-normalized Fourier transforms converge to the
target \(\Xi\) of AC-S33 uniformly on each closed substrip
\(|\operatorname{Im}z|\le b<1/2\).
Condition (LI.5) remains an additional lower-bound hypothesis;
it is not proved here.
\end{theorem}

\begin{proof}
\proofstep{1. Exact size and logarithmic moments.}
Changing variables \(x=\lambda e^t\) in (LI.2) gives exactly
\[
       \|w_n\|_2^2
          =\int_\lambda^\infty |y_n(x)|^2\,dx
          \le \delta_n.
       \tag{LI.7}
\]
This uses the actual concentration defect.  In the even normalization
the integral is half the full two-sided defect, but the displayed
inequality suffices and avoids any change of convention.
The already proved exterior estimate (SE.12) gives
\[
              |w_n(t)|\le C_n\lambda\sqrt{\delta_n}\,e^{-t/2}
                         \quad(t>0).
              \tag{LI.8}
\]
For \(\lambda\ge2\), put \(R=4\log\lambda\).
On \(0<t<R\), use (LI.7).  On \(t\ge R\), use (LI.8).
For each fixed nonnegative integer \(m\), integration by parts
repeatedly, or expansion of \((R+v)^{2m}\), gives
\[
 \int_R^\infty t^{2m}e^{-t}\,dt
                 \le C_m e^{-R}(1+R)^{2m}.
\]
Therefore
\[
 \begin{split}
 \|t^m w_n\|_2^2
 &\le R^{2m}\delta_n+
       C_{n,m}\lambda^2\delta_n e^{-R}(1+R)^{2m}\\
 &\le C_{n,m}\delta_n(1+\log\lambda)^{2m}.
 \end{split}
 \tag{LI.9}
\]
For \(m=0\), the exact estimate (LI.7) is stronger.
No bound on a logarithmic derivative is used in this step.

\proofstep{2. Interpolation preserves the exact leakage scale.}
Apply (SE.2) at the fixed order \(s_0=1/4\):
\[
              \|t^m w_n\|_{H^{1/4}}
                       \le C_{n,m}\lambda^2\sqrt{\delta_n}.
              \tag{LI.10}
\]
In particular the actual jump of \(w_n\) at zero has already been
included; we make no smooth-extension assertion.  For every
\(f\in H^{1/4}(\mathbb R)\) and \(0<s<1/4\), Holder's inequality
on its Fourier transform gives
\[
 \|f\|_{H^s}
       \le \|f\|_2^{1-4s}\|f\|_{H^{1/4}}^{4s}.
 \tag{LI.11}
\]
Indeed write the squared integrand with weight \((1+\xi^2)^s\)
as the product of the \((1-4s)\)-th power of
\(|\widehat f|^2\) and the \(4s\)-th power of
\((1+\xi^2)^{1/4}|\widehat f|^2\), then integrate.
Equations (LI.9)--(LI.11) prove the second part of (LI.3).

\proofstep{3. The original theta trial and zero sampling.}
Let \(w\) be the logarithmic Fourier tail of the original two-mode
trial, exactly as in (SR.11) and (SR.19).  Its two coefficients in
the unit-norm prolate basis are bounded, by AC-S30.  Since
\(\delta_0\le\delta_4=\delta_\lambda\), (LI.3) gives, for every
\(0<\sigma<1/4\),
\[
 \|w\|_{H^\sigma}+\|tw\|_{H^\sigma}
      \le C_\sigma\sqrt{\delta_\lambda}\,
                     \lambda^{8\sigma}(1+\log\lambda).
 \tag{LI.12}
\]
This uses only fixed modes \(0,4\), so there is no growing-index
constant or change to the original mean-zero constraint.

Fix \(0<\sigma<1/4\).  Choose \(0<s<s'<\sigma\) and
\(b_0>0\) with \(s'+b_0<\sigma\).  Under \(A\), the already
proved Mellin multiplier and derivative estimate (SR.17), followed
by sharp half-line restriction (SR.20), bounds the two norms in
the sampling inequality (SR.22) by the right side of (LI.12),
apart from the explicit scalar term in (SR.19).
That scalar satisfies \(|c_Q|\le C\delta_\lambda\), and
\(1_{t>0}e^{-t/2}\) and its first moment lie in \(H^s\) for
\(s<1/2\).  Its contribution is at most
\(C_s\lambda^{-1/2}\delta_\lambda\), which is bounded by the
right side of (LI.12), since \(0<\delta_\lambda\le1\).
Consequently the omitted lower theta tail \(R\) obeys
\[
 \|R\|_{H^s}+\|tR\|_{H^s}
      \le C_\sigma\sqrt{\delta_\lambda}\,
                       \lambda^{8\sigma}(1+\log\lambda).
 \tag{LI.13}
\]
Here \(R\) uses the translated variable of (SR.19).
The subsequent reflection and translation have modulus-one Fourier
factors at the real zero ordinates; no extra logarithmic translation
moment is being inserted into (LI.13).

The exact closed-form identity (SR.23) and its multiplicity-weighted
sampling bound now yield
\[
 q_L(k_\lambda)
       \le C_\sigma\lambda^{16\sigma}
                    (1+\log\lambda)^2\delta_\lambda.
 \tag{LI.14}
\]
Given \(\varepsilon>0\), choose \(16\sigma<\varepsilon\) and
absorb the remaining fixed logarithmic square into
\(\lambda^{\varepsilon-16\sigma}\).  This proves the first part
of (LI.4).  Under \(A\), reflection is a form isometry and the
even and odd form components are nonnegative and orthogonal, so
the even part has no larger form value.  Its norm tends to
\(\|K\|_2>0\), by (LI.1).
The minimum principle and the same classical fixed-index Fuchs
asymptotic (SR.6),
\(\delta_\lambda\asymp\lambda^9e^{-4\pi\lambda^2}\), give
the remaining assertion in (LI.4).

\proofstep{4. The reduced sufficient lower bound.}
Assume (LI.5), and initially choose \(0<\varepsilon<1\).
Let \(v_\lambda\) be the normalized even part of \(k_\lambda\),
with the phase aligned with \(K\).  Then
\[
  q_L(v_\lambda)/a_2(L)
        \le C\lambda^{-(1-\varepsilon)}
                 \longrightarrow0.
  \tag{LI.15}
\]
Since the actual form is nonnegative under \(A\), its minimum is
strictly below \(a_2\) eventually, so its bottom eigenvalue is
simple.  If that line were odd, the even vector \(v_\lambda\)
would be orthogonal to it and have Rayleigh quotient at least
\(a_2\), contradicting (LI.15).  Hence it is even.

Expand \(v_\lambda\) in an orthonormal eigenbasis, and choose the
unit bottom vector \(u_\lambda\) with positive inner product with
\(v_\lambda\).  Nonnegativity and the definition of \(a_2\) give
\[
 1-|\langle u_\lambda,v_\lambda\rangle|^2
                  \le q_L(v_\lambda)/a_2(L),
 \qquad
 \|u_\lambda-v_\lambda\|_2
                  \le C_\varepsilon\lambda^{-(1-\varepsilon)/2}.
 \tag{LI.16}
\]
Both vectors have support in \([-L,L]\).
Cauchy--Schwarz converts the latter bound into a weighted \(L^1\)
bound with factor \(C_b\lambda^b\) for \(b>0\), or
\(\sqrt{2L}\) for \(b=0\).  For each fixed \(b<1/2\), choose \(\varepsilon<1-2b\).
The errors then tend to zero, including the square-root logarithm
when \(b=0\).  The phase of \(u_\lambda\) is fixed by its positive
inner product with \(v_\lambda\), independently of this choice.
The weighted convergence of the original trial
from AC-S13, with AC-S30's scalar, applies also to its even part.
Equation (LI.1) changes its norm from \(\|K\|_2\) by
\(O(\lambda^{-1/2})\); the resulting weighted error also tends
to zero for \(b<1/2\).  This proves (LI.6).  Taking absolute
values under the Fourier integral proves the asserted uniform
convergence on each closed substrip.
\end{proof}

\paragraph{A fixed-strip variant.}
The same proof shows that, for a specified \(0\le b<1/2\),
the weaker lower bound
\[
 a_2(\log\lambda)\ge
           c\lambda^{9+2b+\alpha}e^{-4\pi\lambda^2},
                  \qquad c,\alpha>0,
 \tag{LI.17}
\]
suffices for eventual simplicity and evenness and for (LI.6) at
that particular \(b\).  The ratio in (LI.15) is then
\(O(\lambda^{-(2b+\alpha-\varepsilon)})\), so the weighted
spectral error is \(O(\lambda^{-(\alpha-\varepsilon)/2})\)
when \(b>0\), and has only an additional square-root logarithm
when \(b=0\).  The trial normalization error still requires
\(b<1/2\), as before.

\paragraph{The asymptotic exponent threshold.}
In fact the single condition
\[
 \liminf_{\lambda\to\infty}
       \frac{\log\bigl(a_2(\log\lambda)/\delta_\lambda\bigr)}
            {\log\lambda}\ge1
 \tag{LI.17a}
\]
is sufficient for all the conclusions following (LI.5).
For every \(\eta>0\), it gives
\(a_2(\log\lambda)\ge\delta_\lambda\lambda^{1-\eta}\)
eventually.  By (LI.4), the trial-to-second-eigenvalue ratio is then
at most \(C_\varepsilon\lambda^{-(1-\eta-\varepsilon)}\).
Choosing \(\eta+\varepsilon<1\) proves eventual simplicity and
evenness as before.  For each prescribed \(b<1/2\), choose instead
\(\eta+\varepsilon<1-2b\); the same Cauchy--Schwarz argument proves
the weighted limit.  Equation (LI.5) implies (LI.17a) by (SR.6).
This criterion still requires a lower-bound input, including control
of the full liminf rather than a subsequence.

\paragraph{The multiplicity limitation also persists at the improved scale.}
Under \(A\), if a nontrivial zero at \(\gamma\ne0\) is multiple,
then the same two-dimensional construction in AC-S34 now gives
\[
 a_2(\log\lambda)\le
        C_{\gamma,\varepsilon}\lambda^\varepsilon\delta_\lambda
      \le C'_{\gamma,\varepsilon}
                  \lambda^{9+\varepsilon}e^{-4\pi\lambda^2}
                         \quad(\varepsilon>0).
 \tag{LI.18}
\]
We give the precise improved input to that construction.  Apply
(LI.3) with \(m=0,1,2\), and the derivative version of (SR.17)
used in AC-S34, to obtain
\[
       \|(1+t^2)R\|_2
          \le C_\sigma\lambda^{8\sigma}
                      (1+\log\lambda)^2\sqrt{\delta_\lambda}.
\]
The same Cauchy--Schwarz argument as (SM.2)--(SM.3), and the factor
\(1+L\) from the first Fourier jet after translation, give
\[
       |e_\lambda|
          \le C_{\gamma,\sigma}\lambda^{8\sigma}
                      (1+\log\lambda)^3\sqrt{\delta_\lambda}.
\]
The fixed real even two-jet correctors in (SM.4)--(SM.5) therefore
cost at most \(C_{\gamma,\varepsilon}\lambda^\varepsilon
\delta_\lambda\) in the positive form, after choosing \(\sigma\)
small and absorbing the logarithmic factors.  Together with (LI.4),
this replaces (SM.6) by that same improved bound for \(j_\lambda\).
Its correction is still exponentially small, so the exact compact
division (SM.8)--(SM.11) and the positive limiting Gram bound are
unchanged.  At the removed two ordinates both values still vanish,
and every other denominator is bounded below by the same fixed
\(d_\gamma>0\); hence (SM.12) gives the same improved form bound
for \(g_\lambda\).  Positivity bounds the cross term, and the
two-dimensional minimum principle proves (LI.18), also in the
even sector.  Thus (LI.5), (LI.17a), or any fixed-strip hypothesis (LI.17),
together with \(A\), would imply simplicity of all nontrivial
zeta zeros.  This remains a limitation of these sufficient lower
bounds, not a necessary condition asserted for every possible
ground-line convergence theorem.

\paragraph{Exact scope and source boundary.}
The tail estimates (LI.3) are unconditional.  The theta and zero-sampling
passage uses \(A\) through the already audited arbitrarily small
Sobolev-order zeta multiplier bound; no zero simplicity is assumed.
No new endpoint asymptotic, derivative of a Fuchs formula, WKB
estimate, or numerical certificate is used.  The only classical
prolate asymptotic is the same fixed-index theorem and stated
original-proof access limitation already recorded in AC-S30.
The lower bound (LI.5), (LI.17), or (LI.17a), remains unproved.
The result removes the fixed two-power loss in the previous open
gate (SE.14), and requires no extra positive exponent at the
full-strip threshold \(10\); it does not prove global
simple-even behavior or unconditional convergence of the bottom line.

\paragraph{Analytic and geometric dictionary.}
The object sharpened here is the actual finite-window Weil-form
Hilbert estimate and its spectral-projection consequence.  Its
inputs are Fourier concentration, the theta multiplier and the
zero-sampling form.  This argument constructs no new etale or
motivic cohomology, arithmetic surface, or square construction.
The ground line is selected only after the explicit second-eigenvalue
lower-bound condition; no geometric identification supplies that
condition in this module.
